
\section{Conclusion}
\label{sec:conclusion}

We introduced LGDWT-GS, a frequency-aware extension of 3DGS designed for few-shot 3D reconstruction. By integrating global and patch-wise DWT supervision, our method captures multiscale structural and textural information, mitigating the overreconstruction tendencies of sparse-view setups. The model preserves both large-scale consistency and fine-grained detail, achieving superior stability and perceptual quality compared to strong baselines such as 3DGS, SparseNeRF, and PGDGS across LLFF, MipNeRF360, and our new greenhouse datasets. Beyond the RGB domain, we extended the framework to a multispectral setting that jointly reconstructs RGB and NIR channels under shared geometry, ensuring spectral and spatial coherence across bands. To support research in this direction, we released a controlled multispectral greenhouse dataset and an accompanying few-shot benchmarking package that standardizes sparse-view evaluation protocols.

Overall, LGDWT-GS demonstrates that incorporating frequency-domain priors and multispectral supervision is an effective strategy for constructing efficient, detail-preserving 3D scene representations under data-limited conditions. Future work will focus on extending this framework with frequency-guided densification and pruning strategies to adaptively refine Gaussian distributions based on spectral frequency cues and to generalize the approach to broader multispectral datasets. Additionally, we aim to extend this framework to “in-the-wild” or less controlled agricultural environments, such as field-grown crops, to enable robust real-world 3D reconstruction under challenging outdoor conditions~\cite{zhang2025wheat3dgs}.


